\subsection{Core Set Mutation Operations}

A \texttt{set} is a mutable hash-based collection. Its elements are unique, unordered, and hashable. Mutation means that the set object itself remains the same object, but its internal membership domain changes.

This is different from creating a new set. When \texttt{.add()}, \texttt{.remove()}, \texttt{.discard()}, \texttt{.pop()}, or \texttt{.clear()} is used, Python modifies the existing set object in place.

\begin{verbatim}
names = {"Homer", "Marge"}

print(f"before: {names!r}")
print(f"type(names): {type(names)}")
print(f"id(names): {id(names)}")

names.add("Lisa")

print()
print(f"after:  {names!r}")
print(f"type(names): {type(names)}")
print(f"id(names): {id(names)}")
\end{verbatim}

Possible output:

\begin{verbatim}
before: {'Homer', 'Marge'}
type(names): <class 'set'>
id(names): 2199772358016

after:  {'Lisa', 'Homer', 'Marge'}
type(names): <class 'set'>
id(names): 2199772358016
\end{verbatim}

The memory identity printed by \texttt{id()} remains the same. The set object was not replaced. Its contents were changed.

\subsubsection{Adding Elements with \texttt{.add()}}

The method \texttt{.add()} inserts one element into a set. If an equal element is already present, the set remains unchanged.

\begin{verbatim}
quotes = set()

print(f"start: {quotes!r}")
print(f"type(quotes): {type(quotes)}")

quotes.add("D'oh!")
quotes.add("No soup for you!")
quotes.add("D'oh!")
quotes.add(42)

print()
print(f"after additions: {quotes!r}")
print(f"len(quotes): {len(quotes)}")
\end{verbatim}

Possible output:

\begin{verbatim}
start: set()
type(quotes): <class 'set'>

after additions: {"D'oh!", 42, 'No soup for you!'}
len(quotes): 3
\end{verbatim}

The value \texttt{"D'oh!"} was added twice, but it appears only once. The uniqueness invariant is enforced during insertion.

The method itself returns \texttt{None}. Its purpose is to mutate the set, not to produce a new set.

\begin{verbatim}
names = {"Arthur", "Brian"}

result = names.add("Lancelot")

print(f"names: {names!r}")
print(f"result: {result!r}")
print(f"type(result): {type(result)}")
\end{verbatim}

Possible output:

\begin{verbatim}
names: {'Lancelot', 'Arthur', 'Brian'}
result: None
type(result): <class 'NoneType'>
\end{verbatim}

This is common for in-place mutation methods in Python. The changed object is \texttt{names}; the method call itself does not return the changed set.

The inserted element must be hashable.

\begin{verbatim}
items = set()

items.add("spam")
items.add(("eggs", 42))

print(f"items: {items!r}")

try:
    items.add(["larch", 42])
except TypeError as err:
    print(f"add failed: {err}")
\end{verbatim}

Possible output:

\begin{verbatim}
items: {('eggs', 42), 'spam'}
add failed: unhashable type: 'list'
\end{verbatim}

The tuple is accepted because its elements are hashable. The list is rejected because lists are mutable and unhashable.

We can inspect the set methods relevant to this point.

\begin{verbatim}
items = {"spam"}

print(f"type(items): {type(items)}")
print(f"add in dir(items): {'add' in dir(items)}")
print(f"__contains__ in dir(items): {'__contains__' in dir(items)}")
print(f"__iter__ in dir(items): {'__iter__' in dir(items)}")
print(f"__len__ in dir(items): {'__len__' in dir(items)}")
\end{verbatim}

Possible output:

\begin{verbatim}
type(items): <class 'set'>
add in dir(items): True
__contains__ in dir(items): True
__iter__ in dir(items): True
__len__ in dir(items): True
\end{verbatim}

The important methods are:

\begin{itemize}
    \item \texttt{add}: inserts one hashable element;
    \item \texttt{\_\_contains\_\_}: supports membership testing with \texttt{in};
    \item \texttt{\_\_iter\_\_}: allows iteration over the set's current elements;
    \item \texttt{\_\_len\_\_}: supports \texttt{len()}.
\end{itemize}

\subsubsection{Removing Elements with \texttt{.remove()}, \texttt{.discard()}, \texttt{.pop()}, and \texttt{.clear()}}

A set has several removal operations because different situations need different failure behavior.

The method \texttt{.remove()} removes a specific element. If the element is not present, Python raises \texttt{KeyError}.

\begin{verbatim}
names = {"Homer", "Marge", "Lisa"}

print(f"before: {names!r}")

names.remove("Marge")

print(f"after:  {names!r}")

try:
    names.remove("Bart")
except KeyError as err:
    print(f"remove failed: {err}")
\end{verbatim}

Possible output:

\begin{verbatim}
before: {'Lisa', 'Homer', 'Marge'}
after:  {'Lisa', 'Homer'}
remove failed: 'Bart'
\end{verbatim}

Use \texttt{.remove()} when absence should be treated as an error. The method says: remove this element, and fail loudly if it is not there.

The method \texttt{.discard()} also removes a specific element, but it does not raise an error if the element is absent.

\begin{verbatim}
names = {"Homer", "Marge", "Lisa"}

print(f"before: {names!r}")

names.discard("Marge")
names.discard("Bart")

print(f"after:  {names!r}")
\end{verbatim}

Possible output:

\begin{verbatim}
before: {'Lisa', 'Homer', 'Marge'}
after:  {'Lisa', 'Homer'}
\end{verbatim}

Use \texttt{.discard()} when absence is acceptable. The method says: remove this element if it exists; otherwise do nothing.

Both \texttt{.remove()} and \texttt{.discard()} return \texttt{None}.

\begin{verbatim}
items = {"spam", "eggs", "larch"}

result1 = items.remove("spam")
result2 = items.discard("beans")

print(f"items: {items!r}")
print(f"result1: {result1!r}")
print(f"result2: {result2!r}")
\end{verbatim}

Possible output:

\begin{verbatim}
items: {'eggs', 'larch'}
result1: None
result2: None
\end{verbatim}

The method \texttt{.pop()} removes and returns one arbitrary element from the set.

\begin{verbatim}
items = {"spam", "eggs", "larch"}

print(f"before: {items!r}")

removed = items.pop()

print(f"removed: {removed!r}")
print(f"type(removed): {type(removed)}")
print(f"after:   {items!r}")
\end{verbatim}

Possible output:

\begin{verbatim}
before: {'eggs', 'spam', 'larch'}
removed: 'eggs'
type(removed): <class 'str'>
after:   {'spam', 'larch'}
\end{verbatim}

The word ``arbitrary'' matters. Since a set is not a positional sequence, \texttt{.pop()} does not mean ``remove the first element'' or ``remove the last element.'' Those ideas belong to ordered containers, not to sets.

Calling \texttt{.pop()} on an empty set raises \texttt{KeyError}.

\begin{verbatim}
items = set()

try:
    removed = items.pop()
except KeyError as err:
    print(f"pop failed: {err}")
\end{verbatim}

Possible output:

\begin{verbatim}
pop failed: 'pop from an empty set'
\end{verbatim}

The method \texttt{.clear()} removes all elements from the set.

\begin{verbatim}
items = {"Arthur", "Brian", "Lancelot", 42}

print(f"before: {items!r}")
print(f"len before: {len(items)}")

items.clear()

print()
print(f"after: {items!r}")
print(f"len after: {len(items)}")
print(f"type(items): {type(items)}")
\end{verbatim}

Possible output:

\begin{verbatim}
before: {'Lancelot', 42, 'Arthur', 'Brian'}
len before: 4

after: set()
len after: 0
type(items): <class 'set'>
\end{verbatim}

The set still exists after \texttt{.clear()}. It is simply empty.

The removal methods can be summarized as follows:

\begin{center}
\begin{tabular}{lll}
\textbf{Method} & \textbf{Meaning} & \textbf{Absent Element Behavior} \\
\hline
\texttt{.remove(x)} & Remove \texttt{x} & Raises \texttt{KeyError} \\
\texttt{.discard(x)} & Remove \texttt{x} if present & Does nothing \\
\texttt{.pop()} & Remove and return one arbitrary element & Raises \texttt{KeyError} if empty \\
\texttt{.clear()} & Remove all elements & Always leaves an empty set \\
\end{tabular}
\end{center}

A complete small script shows the operations together.

\begin{verbatim}
words = {"spam", "eggs", "larch", "ni"}

print(f"start: {words!r}")

words.remove("eggs")
print(f"after remove('eggs'):   {words!r}")

words.discard("beans")
print(f"after discard('beans'): {words!r}")

one_word = words.pop()
print(f"popped: {one_word!r}")
print(f"after pop():            {words!r}")

words.clear()
print(f"after clear():          {words!r}")
\end{verbatim}

Possible output:

\begin{verbatim}
start: {'eggs', 'ni', 'spam', 'larch'}
after remove('eggs'):   {'ni', 'spam', 'larch'}
after discard('beans'): {'ni', 'spam', 'larch'}
popped: 'ni'
after pop():            {'spam', 'larch'}
after clear():          set()
\end{verbatim}

\subsubsection{Membership Testing with \texttt{in}}

Membership testing is the central read operation for sets. The expression \texttt{x in some\_set} asks whether \texttt{x} is currently a member of the set.

\begin{verbatim}
allowed = {"Homer", "Marge", "Lisa", "Bart"}

candidate1 = "Lisa"
candidate2 = "Burns"

print(f"allowed: {allowed!r}")
print()

print(f"{candidate1!r} in allowed: {candidate1 in allowed}")
print(f"{candidate2!r} in allowed: {candidate2 in allowed}")
\end{verbatim}

Possible output:

\begin{verbatim}
allowed: {'Lisa', 'Homer', 'Bart', 'Marge'}

'Lisa' in allowed: True
'Burns' in allowed: False
\end{verbatim}

The operator \texttt{in} calls the set's membership machinery. Internally, this depends on hashing and equality, not positional scanning.

The searched value must also be hashable.

\begin{verbatim}
items = {"spam", "eggs", 42}

print(f"{'spam' in items = }")
print(f"{42 in items = }")

try:
    print(["spam"] in items)
except TypeError as err:
    print(f"membership test failed: {err}")
\end{verbatim}

Possible output:

\begin{verbatim}
'spam' in items = True
42 in items = True
membership test failed: unhashable type: 'list'
\end{verbatim}

The list cannot be searched as a set member because checking membership would require computing its hash. Since the list is unhashable, the membership operation fails.

The opposite operator, \texttt{not in}, checks absence.

\begin{verbatim}
blocked = {"Newman", "Burns", "Sideshow Bob"}

name = "Jerry"

print(f"blocked: {blocked!r}")
print(f"name: {name!r}")
print(f"name not in blocked: {name not in blocked}")
\end{verbatim}

Possible output:

\begin{verbatim}
blocked: {'Sideshow Bob', 'Burns', 'Newman'}
name: 'Jerry'
name not in blocked: True
\end{verbatim}

Membership testing is often used before mutation when the program needs explicit control over what happens next.

\begin{verbatim}
names = {"Arthur", "Brian"}

new_name = "Arthur"

if new_name in names:
    print(f"{new_name!r} is already present")
else:
    names.add(new_name)
    print(f"{new_name!r} was added")

print(f"names: {names!r}")
\end{verbatim}

Possible output:

\begin{verbatim}
'Arthur' is already present
names: {'Arthur', 'Brian'}
\end{verbatim}

However, this check is not needed merely to preserve uniqueness. The set already does that.

\begin{verbatim}
names = {"Arthur", "Brian"}

names.add("Arthur")
names.add("Lancelot")

print(f"names: {names!r}")
print(f"len(names): {len(names)}")
\end{verbatim}

Possible output:

\begin{verbatim}
names: {'Lancelot', 'Arthur', 'Brian'}
len(names): 3
\end{verbatim}

The explicit membership check is useful when the program wants a different action depending on presence or absence. It is not required for safe insertion.

A final example combines membership and mutation.

\begin{verbatim}
seen = set()

line1 = "Nobody expects the Spanish Inquisition!"
line2 = "D'oh!"
line3 = "Nobody expects the Spanish Inquisition!"

for line in [line1, line2, line3]:
    if line in seen:
        print(f"duplicate ignored: {line!r}")
    else:
        seen.add(line)
        print(f"new line stored:   {line!r}")

print()
print(f"seen: {seen!r}")
print(f"len(seen): {len(seen)}")
\end{verbatim}

Possible output:

\begin{verbatim}
new line stored:   'Nobody expects the Spanish Inquisition!'
new line stored:   "D'oh!"
duplicate ignored: 'Nobody expects the Spanish Inquisition!'

seen: {'Nobody expects the Spanish Inquisition!', "D'oh!"}
len(seen): 2
\end{verbatim}

The set acts as a memory of already-seen values. It does not remember where they appeared, and it does not remember how many times they appeared. It remembers membership.